\clearpage
\SourceLeaf{412}
\begin{center}\large ZUR TRANSFORMATION DER FLÄCHEN.\end{center}
\noindent\hspace*{1.2em}Die Punkte einer an sich gegebenen Fläche, deren Form noch unbestimmt bleibt, werden durch $r$ und $\theta$ bestimmt. Dabei bezeichnet $r$ den Abstand von einem gegebenen Punkt, gemessen entlang einer kürzesten Linie, und $\theta$ den Winkel dieser kürzesten Linie mit einer festen anderen [kürzesten Linie], die ebenfalls durch diesen Punkt geht. Jede Linie, für die $\theta$ konstant ist, wird dann von jeder anderen, für die $r$ konstant ist, im rechten Winkel geschnitten. Setzt man ein Element der letzteren Linie $=md\theta$, so hängt $m$ nur von der Natur der Fläche ab, jedoch nicht von ihrer Form, und ist somit eine gegebene Funktion von $r$ und $\theta$.

Sobald man der Fläche eine bestimmte Form gegeben hat, gilt für eine Linie, auf der $\theta$ konstant ist, Folgendes.

Es seien $x,y,z$ unbestimmte Koordinaten eines Punkts dieser Linie. Ferner sei
\[
\begin{array}{lll}
\displaystyle\frac{\partial x}{\partial r}=\xi,&\displaystyle\frac{\partial y}{\partial r}=\eta,&\displaystyle\frac{\partial z}{\partial r}=\zeta,\\[6pt]
\displaystyle\frac{\partial\xi}{\partial r}=pX,&\displaystyle\frac{\partial\eta}{\partial r}=pY,&\displaystyle\frac{\partial\zeta}{\partial r}=pZ,
\end{array}
\]
\noindent wobei
\[
p=\BareRoot{\left(\frac{\partial\xi^2}{\partial r^2}+\frac{\partial\eta^2}{\partial r^2}+\frac{\partial\zeta^2}{\partial r^2}\right)},
\]
\noindent also ist $\dfrac1p$ der Krümmungsradius, und es gilt
\[
XX+YY+ZZ=1.
\]
